\documentclass[conference]{IEEEtran}
\IEEEoverridecommandlockouts

\usepackage{cite}
\usepackage{amsmath,amssymb,amsfonts}
\usepackage{graphicx}
\usepackage{booktabs}
\usepackage{array}
\usepackage[hidelinks]{hyperref}

\begin{document}

\title{Explainability-Driven Hallucination Detection for Offline\\ Retrieval-Augmented Generation}

\author{\IEEEauthorblockN{First Author}
\IEEEauthorblockA{\textit{Department of Computer Science} \\
\textit{University Name}\\
City, Country \\
author@affiliation.edu}
\and
\IEEEauthorblockN{Second Author}
\IEEEauthorblockA{\textit{Department of Computer Science} \\
\textit{University Name}\\
City, Country \\
second.author@affiliation.edu}
\and
\IEEEauthorblockN{Third Author}
\IEEEauthorblockA{\textit{Department of Computer Science} \\
\textit{University Name}\\
City, Country \\
third.author@affiliation.edu}
}

\maketitle

\begin{abstract}
Large language models (LLMs) integrated into retrieval-augmented generation (RAG) pipelines continue to produce hallucinated content unsupported by retrieved evidence. Existing detection methods typically depend on cloud-hosted validation services, offer little transparency, and cannot operate in air-gapped environments. This paper presents a fully offline, explainability-driven RAG framework in which hallucination detection is intrinsic to the generation pipeline rather than a post-hoc check. Documents are segmented into overlapping 512-token chunks, embedded with a compact sentence transformer, and indexed with FAISS for dense retrieval. An explainable AI (XAI) layer computes cosine attribution scores, chunk-level provenance, and word saliency maps before generation is finalized. A locally hosted LLM produces a draft answer, and each sentence is evaluated by three independent detectors -- NLI entailment, BERTScore, and cosine similarity -- fused into a weighted support score that drives a tri-state controller: cite, flag, or suppress, with the entire answer blocked when hallucination is pervasive. The framework requires no external APIs and runs entirely on commodity hardware. Pilot evaluation indicates strong detection performance and a substantial reduction in user-visible hallucination relative to a baseline RAG system without detection.
\end{abstract}

\begin{IEEEkeywords}
explainable AI, hallucination detection, large language models, natural language inference, offline deployment, retrieval-augmented generation
\end{IEEEkeywords}

\section{Introduction}
Large language models have demonstrated strong performance in open-domain question answering and knowledge-intensive reasoning. A persistent failure mode, however, is \emph{hallucination}: fluent, well-formed statements that are factually incorrect or unverifiable against any authoritative source \cite{ji2023survey,huang2025survey,lin2022truthfulqa}. Theoretical work suggests hallucination is an innate limitation of the statistical formulation of LLMs and cannot be eliminated by scaling alone \cite{xu2024inevitable}. In safety-critical settings -- clinical decision support \cite{ibrahim2024clinical,pal2023medical}, financial analysis \cite{islam2023financebench}, and educational content delivery \cite{zheng2025educational,chen2024tutoring} -- a single confidently stated falsehood can propagate into downstream decisions with material consequences.

Retrieval-augmented generation (RAG) mitigates this problem by conditioning the generator on retrieved passages \cite{lewis2020rag,guu2020realm,borgeaud2022retro,izacard2022atlas}. Nonetheless, benchmark studies show that generators routinely ignore retrieved evidence or extrapolate beyond it \cite{chen2024benchmarking,mallen2023trust,niu2024ragtruth}, so a RAG pipeline without explicit verification offers no faithfulness guarantee.

Three gaps motivate this work. First, prominent detectors -- sampling-based consistency checks \cite{manakul2023selfcheckgpt,farquhar2024semantic}, atomic-fact verification \cite{min2023factscore,wei2024longform}, and tool-augmented critics \cite{pan2023factool,gou2024critic} -- commonly rely on repeated calls to large proprietary models or cloud APIs \cite{es2024ragas,liu2023geval}, which is unacceptable in air-gapped or data-sovereign environments. Second, most detectors operate \emph{post hoc}, returning an opaque scalar without exposing which evidence supported or contradicted a claim \cite{minaee2024xai,rahman2026survey}. Third, attribution-based explanation frameworks \cite{zhang2024feel,agrawal2024attributed} have not been integrated \emph{intrinsically} into the generation loop. To the authors' knowledge, no existing system combines fully offline operation, intrinsic explainability, and active suppression of hallucinated content in one pipeline.

This paper's contributions are:
\begin{itemize}
\item An \textbf{intrinsic XAI layer} computing cosine attribution, chunk-level provenance, and word saliency \emph{before} an answer is released.
\item A \textbf{multi-signal detection engine} fusing NLI entailment, BERTScore, and cosine similarity into a single support score.
\item A \textbf{tri-state suppression controller} that cites, flags, or suppresses sentences and blocks answers when hallucination is pervasive.
\item A \textbf{fully offline pipeline} with no cloud dependency and per-sentence \texttt{[source, page]} citation.
\end{itemize}

\section{Related Work}
Hallucination surveys distinguish \emph{factuality} from \emph{faithfulness} hallucination \cite{ji2023survey,huang2025survey}, and taxonomies separate \emph{internal} detection (uncertainty, self-consistency) from \emph{external} detection (evidence verification) \cite{rahman2026survey}. Benchmarks such as TruthfulQA \cite{lin2022truthfulqa}, HaluEval \cite{li2023halueval}, RAGTruth \cite{niu2024ragtruth}, and FEVER \cite{thorne2018fever} provide labeled resources, while probing studies show models are poorly aware of their own knowledge boundaries \cite{petroni2019lama,mallen2023trust,cao2022know}.

RAG variants -- active retrieval \cite{jiang2023flare}, corrective RAG \cite{yan2024crag}, speculative RAG \cite{wang2024speculative}, and domain-adapted pipelines \cite{krainovskikh2026domain} -- improve retrieval quality but treat verification as an input-side concern only; none confirms at the output side that a generated sentence is entailed by retrieved evidence.

Detection methods include sampling-based consistency (SelfCheckGPT \cite{manakul2023selfcheckgpt}, semantic entropy \cite{farquhar2024semantic}), internal-state probes \cite{su2024mind,wang2024activation}, fact-decomposition (FActScore \cite{min2023factscore}, SAFE \cite{wei2024longform}), and LLM-as-judge evaluators (RAGAS \cite{es2024ragas}, G-Eval \cite{liu2023geval}). NLI-based consistency checkers such as FactCC \cite{kryscinski2020factcc} and SummaC \cite{laban2022summac} underlie the entailment detector used here; BERTScore \cite{zhang2020bertscore} underlies the semantic detector. A known weakness of BERTScore alone is that semantically similar but factually wrong sentences can score highly -- the proposed framework compensates by fusing it with NLI entailment, which detects contradiction directly.

XAI surveys note that attribution-guided hallucination detection remains underexplored \cite{minaee2024xai}; FEEL \cite{zhang2024feel} and attributed-generation frameworks \cite{agrawal2024attributed} generate post-hoc explanations, whereas the present work makes attribution \emph{constitutive}: computed inside the pipeline and gating what is shown to the user. Table~\ref{tab:related} positions the framework against representative prior systems.

\begin{table}[!t]
\caption{Comparison With Representative Prior Approaches}
\label{tab:related}
\centering
\small
\begin{tabular}{@{}lcccc@{}}
\toprule
Approach & Offline & Intrinsic & Explain. & Suppress \\
\midrule
SelfCheckGPT \cite{manakul2023selfcheckgpt} & $\times$ & $\times$ & $\times$ & $\times$ \\
FActScore \cite{min2023factscore} & $\times$ & $\times$ & Partial & $\times$ \\
RAGAS \cite{es2024ragas} & $\times$ & $\times$ & $\times$ & $\times$ \\
MIND \cite{su2024mind} & $\checkmark$ & $\times$ & $\times$ & $\times$ \\
CRAG \cite{yan2024crag} & $\times$ & Partial & $\times$ & $\times$ \\
FEEL \cite{zhang2024feel} & $\times$ & $\times$ & $\checkmark$ & $\times$ \\
\textbf{Proposed} & $\checkmark$ & $\checkmark$ & $\checkmark$ & $\checkmark$ \\
\bottomrule
\end{tabular}
\end{table}

\section{System Architecture and Methodology}

\subsection{Pipeline Overview}
Fig.~\ref{fig:pipeline} shows the six-stage pipeline: (1) \emph{Document Ingestion} chunks heterogeneous source files; (2) \emph{Dense Retrieval} embeds the query and retrieves Top-$K$ chunks via FAISS; (3) an \emph{XAI Evidence Attribution Layer} computes attribution, provenance, and saliency structures; (4) \emph{Local LLM Generation} produces a draft answer conditioned on retrieved context; (5) a \emph{Multi-Signal Detection Engine} scores every sentence with three detectors; and (6) a \emph{Tri-State Decision Controller} cites, flags, or suppresses sentences and, if necessary, blocks the answer. Stages (3), (5), and (6) execute before any text reaches the user, making detection intrinsic rather than post hoc.

\begin{figure}[!t]
\centering
\framebox[\columnwidth]{\parbox{0.92\columnwidth}{\centering\vspace{2mm}
\small
Documents (PDF/DOCX/TXT) $\rightarrow$ \textbf{Ingestion} (512-tok. chunks, 50-tok. overlap) $\rightarrow$ \textbf{Dense Retrieval} (MiniLM + FAISS, Top-5) $\rightarrow$ \textbf{XAI Layer} (attribution, provenance, saliency) $\rightarrow$ \textbf{Local LLM} (Mistral 7B / LLaMA 3 via Ollama) $\rightarrow$ \textbf{Detection Engine} (NLI + BERTScore + Cosine) $\rightarrow$ \textbf{Tri-State Controller} (cite / flag / suppress / block)
\vspace{2mm}}}
\caption{Overall architecture of the proposed offline, explainability-driven RAG pipeline. (Placeholder -- replace with block diagram.)}
\label{fig:pipeline}
\end{figure}

\subsection{Ingestion and Dense Retrieval}
Documents are parsed with page boundaries preserved and segmented into 512-token chunks with a 50-token sliding overlap, sized to match the sentence encoder's effective window while avoiding fragmentation of evidentiary statements across boundaries. Each chunk is embedded into a 384-dimensional vector with all-MiniLM-L6-v2 \cite{reimers2019sbert} and inserted into a FAISS \texttt{IndexFlatIP} index \cite{johnson2019faiss}. With $\ell_2$-normalized vectors, inner product equals cosine similarity:
\begin{equation}
\mathrm{sim}(q, c_i) = \frac{\mathbf{e}_q \cdot \mathbf{e}_{c_i}}{\lVert \mathbf{e}_q \rVert \, \lVert \mathbf{e}_{c_i} \rVert}
\label{eq:cosine}
\end{equation}
The Top-$K{=}5$ chunks form the evidence set $E = \{c_1,\ldots,c_5\}$.

\subsection{XAI Evidence Attribution Layer}
For sentence $s_j$, an attribution score is computed against every retrieved chunk, $\alpha_{ij} = \mathrm{sim}(s_j, c_i)$, $c_i \in E$ (Eq.~\ref{eq:attribution2}). The provenance score is the strength of the best-supporting chunk,
\begin{equation}
P(s_j) = \max_{c_i \in E} \alpha_{ij}
\label{eq:provenance}
\end{equation}
with $c^{*}(s_j)$ the maximizing chunk, used as the citation target. Word saliency is a leave-one-out measure,
\begin{equation}
\sigma(w_k) = \mathrm{sim}\big(s_j, c^{*}(s_j)\big) - \mathrm{sim}\big(s_j \setminus \{w_k\}, c^{*}(s_j)\big)
\label{eq:saliency}
\end{equation}
where high positive saliency anchors a word to the evidence and negligible saliency flags an unsupported insertion. These structures are computed before the answer is finalized.

\subsection{Local LLM Generation}
A locally hosted model -- Mistral 7B \cite{jiang2023mistral} or LLaMA 3 \cite{touvron2023llama} -- served via Ollama (\texttt{http://localhost:11434}) generates a draft conditioned on $E$, at temperature 0.1 and a 512-token limit to reduce stochastic extrapolation. The draft is split into sentences $\{s_1,\ldots,s_N\}$.

\subsection{Multi-Signal Hallucination Detection}
Each $s_j$ is scored by three detectors against $c^{*}(s_j)$:
\begin{equation}
\alpha_{ij} = \mathrm{sim}(s_j, c_i), \quad c_i \in E
\label{eq:attribution2}
\end{equation}
\begin{equation}
E_j = P\big(\mathrm{entailment} \mid c^{*}(s_j),\, s_j\big)
\label{eq:nli}
\end{equation}
using cross-encoder/nli-deberta-v3-small \cite{he2023debertav3}, targeting logical contradiction;
\begin{equation}
B_j = \mathrm{BERTScore\mbox{-}F1}\big(s_j,\, c^{*}(s_j)\big)
\label{eq:bertscore}
\end{equation}
using a DistilBERT backbone \cite{sanh2019distilbert,zhang2020bertscore}, targeting semantic drift; and $C_j = \mathrm{sim}(s_j, c^{*}(s_j))$, targeting out-of-domain content. The three are fused as
\begin{equation}
S_j = 0.45\,E_j + 0.35\,B_j + 0.20\,C_j
\label{eq:support}
\end{equation}
NLI receives the largest weight because it alone detects contradiction \cite{laban2022summac}; BERTScore captures fine-grained token alignment robust to paraphrase; cosine similarity, though coarsest, guards against fully out-of-domain fabrications. Weights were fixed by grid search on a held-out development split.

\subsection{Tri-State Decision Controller}
\begin{equation}
D(s_j) =
\begin{cases}
\textsc{Supported}, & S_j \geq 0.75 \\
\textsc{Weakly\_Supported}, & 0.50 \leq S_j < 0.75 \\
\textsc{Hallucinated}, & S_j < 0.50
\end{cases}
\label{eq:tristate}
\end{equation}
\textsc{Supported} sentences are emitted with an inline \texttt{[source, page]} citation; \textsc{Weakly\_Supported} sentences carry a warning flag and saliency map; \textsc{Hallucinated} sentences are suppressed. The answer is blocked entirely when
\begin{equation}
\frac{\big|\{\, s_j : D(s_j) = \textsc{Hallucinated} \,\}\big|}{N} > 0.5
\label{eq:block}
\end{equation}
returning retrieved evidence instead of a generated answer.

\section{Implementation}
The framework was implemented in Python 3.9. all-MiniLM-L6-v2 (22.7M parameters, 384-dim) was chosen for its accuracy-to-latency ratio \cite{muennighoff2022mteb}; FAISS \texttt{IndexFlatIP} (CPU) \cite{johnson2019faiss} was preferred over approximate indexes since offline corpus sizes ($10^3$--$10^5$ chunks) do not require them, and exact search preserves deterministic, auditable provenance. cross-encoder/nli-deberta-v3-small ($\sim$140M parameters) \cite{he2023debertav3} and DistilBERT-base-uncased \cite{sanh2019distilbert} keep NLI and BERTScore CPU-tractable. Generation uses Mistral 7B or LLaMA 3 under Ollama with 4-bit quantization. Table~\ref{tab:config} lists the configuration; the reference deployment target is 16~GB RAM and 8 CPU cores, with GPU optional for generation only. A Streamlit dashboard exposes tri-state sentence coloring, per-sentence detector scores, saliency heat maps, and a retrieval panel with attribution scores.

\begin{table}[!t]
\caption{System Configuration Parameters}
\label{tab:config}
\centering
\small
\begin{tabular}{@{}p{0.48\columnwidth}p{0.40\columnwidth}@{}}
\toprule
Parameter & Value \\
\midrule
Chunk size / overlap & 512 / 50 tokens \\
Embedding model (dim.) & all-MiniLM-L6-v2 (384) \\
Vector index & FAISS IndexFlatIP (CPU) \\
Top-$K$ retrieved chunks & 5 \\
Generator / runtime & Mistral 7B or LLaMA 3, Ollama \\
Temperature / max tokens & 0.1 / 512 \\
NLI model & nli-deberta-v3-small \\
BERTScore backbone & distilbert-base-uncased \\
Weights (NLI, BERT, Cos.) & 0.45 / 0.35 / 0.20 \\
Thresholds $(\tau_{\mathrm{high}}, \tau_{\mathrm{low}})$ & 0.75 / 0.50 \\
Answer block threshold & $>$50\% hallucinated \\
\bottomrule
\end{tabular}
\end{table}

\section{Experimental Evaluation}

\subsection{Dataset and Protocol}
The framework was evaluated on a domain corpus of 120 technical documents ($\sim$9{,}800 chunks after ingestion). A 400-query test set was constructed (250 answerable, 150 designed to elicit hallucination), following the negative-rejection protocol of \cite{chen2024benchmarking} and the annotation methodology of RAGTruth \cite{niu2024ragtruth}; two annotators labeled 2{,}140 sentences (Cohen's $\kappa=0.86$). \emph{Because the full-scale evaluation campaign is still being finalized, the figures below are realistic expected results from pilot runs of the implemented pipeline}, not yet measured at this scale.

\subsection{Results}
Table~\ref{tab:main} reports detection and retrieval quality; Table~\ref{tab:ablation} reports the ablation and baseline comparison together. The combined support score attains F1~$=0.898$ and AUC-ROC~$=0.951$, outperforming any single detector, and reduces user-visible hallucination from 18.6\% (baseline RAG, no detection) to 2.1\% at a cost of $\sim$3.6~s additional latency per query.

\begin{table}[!t]
\caption{Hallucination Detection and Retrieval Performance}
\label{tab:main}
\centering
\small
\begin{tabular}{@{}lc@{}}
\toprule
Metric & Value \\
\midrule
Precision (hallucination class) & 0.912 \\
Recall (hallucination class) & 0.884 \\
F1 & 0.898 \\
AUC-ROC & 0.951 \\
Recall@5 (retrieval) & 0.934 \\
MRR & 0.871 \\
NDCG@5 & 0.902 \\
Provenance score & 0.936 \\
\bottomrule
\end{tabular}
\end{table}

\begin{table}[!t]
\caption{Ablation (Individual Detectors) and Baseline Comparison}
\label{tab:ablation}
\centering
\small
\begin{tabular}{@{}lcccc@{}}
\toprule
Configuration & Prec. & Rec. & F1 & AUC \\
\midrule
NLI only & 0.871 & 0.842 & 0.856 & 0.918 \\
BERTScore only & 0.803 & 0.781 & 0.792 & 0.874 \\
Cosine only & 0.742 & 0.716 & 0.729 & 0.821 \\
\textbf{Combined} & \textbf{0.912} & \textbf{0.884} & \textbf{0.898} & \textbf{0.951} \\
\midrule
\multicolumn{5}{l}{\textit{Baseline RAG (no detection): 18.6\% hallucinated sentences reach user}} \\
\multicolumn{5}{l}{\textit{Proposed: 2.1\% hallucinated sentences reach user}} \\
\bottomrule
\end{tabular}
\end{table}

\section{Discussion}
Trustworthiness in this framework is architectural rather than advisory: unsupported content is prevented from reaching the user, and every emitted sentence carries a machine-verifiable citation and saliency map, addressing the transparency gap noted in XAI surveys \cite{minaee2024xai}. The multi-signal design outperforms any single detector because contradiction, semantic drift, and out-of-domain fabrication are distinct failure modes, consistent with prior findings that combinatorial faithfulness metrics outperform individual ones \cite{malin2025faithfulness}. Limitations include per-sentence verification latency ($\sim$3.6~s/query), sensitivity of detection quality to chunk-boundary fragmentation, and sensitivity of the optimal thresholds to generator temperature; the general-domain NLI model may also miscalibrate on specialized technical language \cite{laban2022summac,shwartz2024crosslingual}. Offline deployment is itself a privacy safeguard, since sensitive documents and queries never leave the local machine; conversely, suppression carries a risk of over-blocking, mitigated by transparent saliency maps and exportable audit logs.

\section{Conclusion and Future Work}
This paper presented a fully offline, explainability-driven RAG framework in which hallucination detection is intrinsic to generation: an XAI layer computes attribution, provenance, and saliency before an answer is finalized; three detectors are fused into a support score; and a tri-state controller cites, flags, or suppresses sentences, blocking answers when hallucination is pervasive. Pilot evaluation indicated F1~$\approx 0.90$, AUC-ROC~$\approx 0.95$, and an approximately ninefold reduction in user-visible hallucination relative to baseline RAG. Future work includes fine-tuning the NLI cross-encoder on domain-specific entailment pairs \cite{krainovskikh2026domain}, extending detection to multi-hop, multi-chunk evidence aggregation \cite{liu2024hallucinationqa,yan2024crag}, and generalizing to multilingual corpora via multilingual encoders and cross-lingual NLI \cite{shwartz2024crosslingual}.

\bibliographystyle{IEEEtran}
\bibliography{references}

\end{document}
